% ============================================================================
% Fichier  : partie1/P01_C01_risque_incertitude.tex
% Titre    : Risque, Incertitude et Décision
% Auteur   : MINKA MI NGUIDJOÏ Thierry Emmanuel
% Version  : 1.0 -- novembre 2026
% ============================================================================
% Sources : ISO 31000:2018 ; ISO Guide 73:2009 ; Knight (1921) ; Taleb
% (2007) ; Perrow (1984). Chiffres de LiZbiZ : cas fictif, cohérents entre eux.
% ============================================================================

\chapter{Risque, Incertitude et Décision}
\label{chap:risque}
\chaptermark{Risque, incertitude, décision}

\epigraph{%
    « Une bonne décision peut mal tourner,\\
    et une mauvaise décision peut bien tourner. »%
}{Principe du cours}

\niveauChapitre{Fondamental}{Prologue~; notions de probabilité.}

\begin{boxobjectifs}
\begin{itemize}
    \item Définir le risque selon l'ISO~31000 et le distinguer de
          l'incertitude au sens de Knight.
    \item Décrire un risque par une distribution de conséquences, et non par
          un seul chiffre.
    \item Calculer et interpréter une perte annuelle attendue.
    \item Juger une décision sur l'information disponible, et non sur son
          résultat.
    \item Reconnaître les limites de la prévision face aux événements rares.
\end{itemize}
\end{boxobjectifs}

\bigskip

% ============================================================================
\section*{Question directrice : le risque de rançongiciel était-il connu en février~?}
% ============================================================================

En février 2026, M.~KOUNOU avait proposé d'exiger un second facteur sur la
passerelle d'accès distant, pour 6~millions de FCFA par an. La dépense
avait été reportée~: « nous n'avons jamais été attaqués ». Le risque
était-il connu~? Pouvait-on le chiffrer~? Et la décision de reporter
était-elle mauvaise, ou seulement malchanceuse~?

\medskip
\textbf{Plan du chapitre.}
\begin{enumerate}
    \item Risque et incertitude (section~\ref{sec:c01-definir}).
    \item Probabilité et conséquences (section~\ref{sec:c01-distribution}).
    \item La perte annuelle attendue (section~\ref{sec:c01-esperance}).
    \item Juger une décision (section~\ref{sec:c01-decision}).
    \item Les événements rares (section~\ref{sec:c01-rares}).
\end{enumerate}

\separateur

% ============================================================================
\section{Risque et Incertitude}
\label{sec:c01-definir}
% ============================================================================

\begin{boxdef}
\termecle{Risque} (ISO~31000)~: effet de l'incertitude sur les objectifs.
Un effet est un écart par rapport à ce qui est attendu~; il peut être
positif ou négatif. Le risque s'exprime souvent par des sources de risque,
des événements, leurs conséquences et leur vraisemblance.
\end{boxdef}

La définition de l'\norme{ISO~31000} et du \norme{Guide~ISO~73} rattache
le risque aux \emph{objectifs}~: sans objectif, pas de risque. Pour
LiZbiZ, l'objectif « facturer chaque jour ouvré » donne son sens au risque
« arrêt de la facturation »~\cite{iso31000_2018,isoguide73}.

En 1921, l'économiste Frank Knight distingue le \termecle{risque}, dont on
connaît la distribution des issues, et l'\termecle{incertitude}, dont on ne
la connaît pas~\cite{knight1921}. Le lancer d'un dé relève du risque~; la
date de la prochaine attaque contre LiZbiZ relève, en grande partie, de
l'incertitude. Le travail du gestionnaire de risques consiste à faire
passer ce qui peut l'être de l'incertitude au risque, par l'observation,
les données et le jugement d'experts, et à reconnaître ce qui n'y passera
pas.

\begin{tableaucours}{L{3.0cm}L{4.6cm}L{4.6cm}}{%
    \tableauHeader{Situation} &
    \tableauHeader{Ce que l'on sait} &
    \tableauHeader{Exemple chez LiZbiZ}
}
Risque & Les issues et leurs probabilités & Pannes de disques, connues
par les statistiques du fabricant \\
Incertitude mesurable & Les issues, et une estimation des probabilités &
Hameçonnage, estimé par les signalements et les études publiées \\
Incertitude profonde & Ni la liste des issues, ni leurs probabilités &
Une rupture technologique, une crise politique régionale \\
\end{tableaucours}
\vspace{-0.8em}
\begin{center}{\small\itshape Tableau~1.1. Du risque à l'incertitude profonde}\end{center}

% ============================================================================
\section{Probabilité et Conséquences}
\label{sec:c01-distribution}
% ============================================================================

Un risque ne se résume pas à un nombre. Il se décrit par la
\termecle{distribution des conséquences}~: pour chaque perte possible, la
probabilité qu'elle survienne sur une période. Deux risques de même perte
moyenne peuvent être très différents.

\begin{tableaucours}{L{3.6cm}L{4.2cm}L{4.2cm}}{%
    \tableauHeader{} &
    \tableauHeader{Risque A~: pannes fréquentes} &
    \tableauHeader{Risque B~: rançongiciel}
}
Fréquence annuelle & 10 pannes & 0,25 événement \\
Perte par événement & 3~millions de FCFA & 120~millions de FCFA \\
Perte annuelle moyenne & 30~millions & 30~millions \\
Pire année plausible & 45~millions & 300~millions ou plus \\
\end{tableaucours}
\vspace{-0.8em}
\begin{center}{\small\itshape Tableau~1.2. Même moyenne, deux risques différents}\end{center}

\textbf{Analyse.} Les deux risques coûtent 30~millions par an en moyenne,
mais le risque~B peut, une année donnée, menacer la trésorerie. Une
entreprise peut absorber A dans son budget~; elle doit transférer ou
réduire B. C'est la \emph{queue} de la distribution, et non sa moyenne,
qui décide.

% ============================================================================
\section{La Perte Annuelle Attendue}
\label{sec:c01-esperance}
% ============================================================================

\begin{boxdef}
\termecle{Perte annuelle attendue}~: espérance de la perte sur un an,
égale, pour un scénario simple, au produit de la fréquence annuelle et de
la perte moyenne par événement. Elle sert à comparer des options, non à
prévoir une année particulière.
\end{boxdef}

\begin{boxscenario}{Opération GAROUA. La décision de février}
\textbf{Hypothèses} (établies après coup, confiance faible à moyenne)~:
fréquence annuelle d'un rançongiciel réussi sans second facteur, 0,25~;
avec second facteur, 0,05~; perte par événement, 300~millions de FCFA
(Tableau~0.4)~; coût du second facteur, 6~millions par an.

\medskip
\textbf{Calcul.} Sans second facteur~: $0{,}25 \times 300 = 75$~millions
par an. Avec~: $0{,}05 \times 300 + 6 = 21$~millions par an.

\medskip
\textbf{Analyse} (confiance moyenne). Même si les fréquences sont
incertaines, l'écart est si grand que la décision ne change pas tant que
la fréquence sans second facteur dépasse $6/300 + 0{,}05 = 0{,}07$ par an.
Reporter la dépense n'était pas justifié par les chiffres~; il était
justifié par l'absence de chiffres.
\end{boxscenario}

Ce raisonnement est un \termecle{seuil de rentabilité}~: au lieu de
chercher la « vraie » fréquence, on cherche la fréquence à partir de
laquelle la décision bascule, puis on se demande si l'on y croit. C'est
souvent plus facile que d'estimer la fréquence elle-même.

% ============================================================================
\section{Juger une Décision}
\label{sec:c01-decision}
% ============================================================================

Une décision se juge sur l'information disponible au moment de la
prendre, et non sur son résultat. Si LiZbiZ avait installé le second
facteur et avait été attaqué par une autre voie, la décision aurait
été bonne et le résultat mauvais. À l'inverse, reporter la dépense et ne
pas être attaqué aurait été une mauvaise décision au résultat favorable.

\begin{tableaucours}{L{3.2cm}L{4.5cm}L{4.5cm}}{%
    \tableauHeader{} &
    \tableauHeader{Résultat favorable} &
    \tableauHeader{Résultat défavorable}
}
Bonne décision & Mérité & Malchance \\
Mauvaise décision & Chance & Conséquence prévisible \\
\end{tableaucours}
\vspace{-0.8em}
\begin{center}{\small\itshape Tableau~1.3. Décision et résultat}\end{center}

\begin{boxalert}
\textbf{Violation typique.} Juger une décision à son résultat (« on n'a
jamais été attaqués, donc on avait raison »).

\smallskip
\textbf{Règle~:} documenter, au moment de décider, les options, les
hypothèses et l'information disponible. Le retour d'expérience juge
ensuite la qualité de ce dossier, et non la chance
(Chapitre~\ref{chap:retex}).
\end{boxalert}

Trois critères éclairent une décision sous incertitude~: la perte attendue
(ce que l'on perd en moyenne), la perte dans un scénario sévère (ce que
l'on ne peut pas se permettre) et la \termecle{valeur de l'information}
(ce que coûterait d'en savoir plus avant de décider). Un dirigeant averse
au risque peut préférer une option plus chère en moyenne si elle supprime
le scénario qui le ruinerait.

% ============================================================================
\section{Les Événements Rares}
\label{sec:c01-rares}
% ============================================================================

Certains événements sont si rares qu'aucune statistique ne permet de les
estimer, et leurs conséquences si grandes qu'ils dominent toutes les
autres. Nassim Taleb les appelle des « cygnes noirs »~: rares, d'impact
extrême, et expliqués après coup comme s'ils avaient été
prévisibles~\cite{taleb2007}. Charles Perrow avait montré, dès 1984, que
dans les systèmes complexes et étroitement couplés, des défaillances
mineures peuvent se combiner en accident majeur~: il les nommait
« accidents normaux »~\cite{perrow1984}.

\begin{boxinfo}
\textbf{Ce que l'on peut faire face à l'imprévisible.} On ne prévoit pas
l'événement, mais on réduit la fragilité~: couplages moins étroits
(segmentation, liaisons coupables entre sites), marges (trésorerie,
stocks), capacité de réaction (cellule de crise entraînée). C'est l'objet
des Parties~IV et~V.
\end{boxinfo}

\begin{boxobjet}{R-01 · Arrêt de la facturation}
\textbf{Scénario~:} un incident informatique arrête la facturation plus
de deux jours ouvrés. \textbf{Cotation~:} fréquence estimée de 0,3 par an,
perte par événement de 40 à 300~millions de FCFA. \textbf{Traitement~:}
réduire (second facteur, sauvegardes hors ligne) et transférer (assurance
cyber). \textbf{Propriétaire~:} Mme~MBALLA. \textbf{Indicateur~:} durée du
dernier test de restauration de la facturation.
\end{boxobjet}

\begin{boxmetier}{Gestionnaire de risques}
\textbf{Chiffrer le seuil, pas seulement le risque.} Quand la fréquence
est incertaine, le gestionnaire présente à la direction la fréquence à
partir de laquelle l'investissement se justifie, et les raisons de croire
qu'elle est dépassée.
\end{boxmetier}

\separateur

\begin{boxresume}
\textbf{Ce qu'il faut retenir de ce chapitre~:}
\begin{itemize}
    \item \textbf{Le risque} est l'effet de l'incertitude sur les
          objectifs (ISO~31000)~; Knight distingue risque et incertitude.
    \item \textbf{Un risque} se décrit par une distribution~; la queue
          compte autant que la moyenne.
    \item \textbf{La perte annuelle attendue} compare des options~; le
          seuil de rentabilité évite d'estimer ce qu'on ne sait pas.
    \item \textbf{Une décision} se juge sur l'information disponible, pas
          sur son résultat.
    \item \textbf{Face aux événements rares}, on réduit la fragilité.
\end{itemize}
\end{boxresume}

\begin{boxinfo}
\textbf{Points de contrôle de l'élève.} Je sais~: définir le risque selon
l'ISO~31000~; distinguer risque et incertitude~; comparer deux risques de
même moyenne~; calculer une perte annuelle attendue et un seuil de
rentabilité~; distinguer bonne décision et bon résultat.
\end{boxinfo}

\bigskip

\begin{boxexo}[title={\ding{45}\quad Exercice~1.1. Définir \niveaubase}]
Donnez, pour LiZbiZ, un objectif, un risque qui s'y rattache, sa source,
l'événement redouté et sa conséquence.
\end{boxexo}

\begin{boxexo}[title={\ding{45}\quad Exercice~1.2. Knight \niveaubase}]
Classez en risque, incertitude mesurable ou incertitude profonde~: la
panne d'un onduleur~; une fraude au virement~; une nouvelle loi sur les
données~; une inondation à Bonabéri.
\end{boxexo}

\begin{boxexo}[title={\ding{45}\quad Exercice~1.3. Calculer \niveaubase}]
Une fraude au virement survient en moyenne une fois tous les quatre ans
pour une perte de 60~millions de FCFA. Quelle est la perte annuelle
attendue~?
\end{boxexo}

\begin{boxexo}[title={\ding{45}\quad Exercice~1.4. Le seuil \niveaumoyen}]
Une double validation des virements coûte 4~millions par an et divise la
fréquence par quatre. À partir de quelle fréquence initiale est-elle
rentable, pour une perte de 60~millions par événement~?
\end{boxexo}

\begin{boxexo}[title={\ding{45}\quad Exercice~1.5. La queue \niveaumoyen}]
Expliquez à Mme~MBALLA pourquoi le risque~B du Tableau~1.2 justifie une
assurance, alors que le risque~A n'en justifie pas.
\end{boxexo}

\begin{boxexo}[title={\ding{45}\quad Exercice~1.6. Décision ou chance \niveaumoyen}]
Classez dans le Tableau~1.3~: LiZbiZ installe le second facteur et n'est
pas attaquée~; elle reporte la dépense et est attaquée~; elle reporte et
n'est pas attaquée.
\end{boxexo}

\begin{boxexo}[title={\ding{45}\quad Exercice~1.7. Le dossier \niveaumoyen}]
Rédigez la note de décision que M.~KOUNOU aurait dû remettre en février~:
options, hypothèses, chiffres, recommandation.
\end{boxexo}

\begin{boxexo}[title={\ding{45}\quad Exercice~1.8. Couplage \niveaumoyen}]
En quoi les liaisons ouvertes entre les sites de LiZbiZ illustrent-elles
le couplage étroit au sens de Perrow~?
\end{boxexo}

\begin{boxexo}[title={\ding{45}\quad Exercice~1.9. Sensibilité \niveauexpert}]
Refaites le calcul de la décision de février avec une perte par
événement de 120~millions au lieu de 300. La conclusion change-t-elle~?
Pour quelle perte bascule-t-elle~?
\end{boxexo}

\begin{boxexo}[title={\ding{45}\quad Exercice~1.10. Fragilité \niveauexpert}]
Proposez cinq mesures qui réduisent la fragilité de LiZbiZ face à un
événement que l'on ne sait pas prévoir. Pour chacune, dites quel couplage
ou quelle marge elle crée.
\end{boxexo}

\bigskip

\begin{boxinfo}
\textbf{Pour aller plus loin.}
\begin{itemize}
    \item ISO~31000~\cite{iso31000_2018}~; Guide~ISO~73~\cite{isoguide73}.
    \item F.~Knight, \emph{Risk, Uncertainty and Profit}~\cite{knight1921}.
    \item C.~Perrow, \emph{Normal Accidents}~\cite{perrow1984}.
\end{itemize}
\end{boxinfo}

\begin{boxinfo}
\textbf{Transition.} Les chiffres ne suffisent pas~: les personnes perçoivent
les risques à travers des biais. Le Chapitre~\ref{chap:perception} les
présente.
\end{boxinfo}
